Работа посвящена выбору результата кластеризации при доступности меток лишь небольшой части объектов. В этой ситуации внешний критерий чувствителен к составу размеченной выборки, а внутренний индекс описывает геометрию данных, которая может не совпадать со структурой целевых классов. Исследуется, при каких условиях внутренний индекс улучшает выбор по ограниченной разметке и когда его использование приводит к потере качества. Рассматривается фиксированный пул разбиений и правило, которое допускает альтернативы лидеру по частичному скорректированному индексу Рэнда с учётом бутстрепной изменчивости парных различий, после чего выбирает кандидата по внутреннему индексу. Для оценки такого вмешательства задаётся протокол сравнения на скрытых проверочных метках с анализом внешнего качества, потерь относительно лучшего кандидата и последствий смены выбранного разбиения. Основные исследуемые факторы --- объём разметки, покрытие классов, согласованность геометрии с целевыми метками и состав пула кандидатов. Цель работы состоит в определении границ применимости совместного использования внутренних и внешних оценок, включая условия систематического ухудшения выбора.
